\begin{savequote}[8cm]
\textlatin{Neque porro quisquam est qui dolorem ipsum quia dolor sit amet, consectetur, adipisci velit...}

There is no one who loves pain itself, who seeks after it and wants to have it, simply because it is pain...
  \qauthor{--- Cicero's \textit{de Finibus Bonorum et Malorum}}
\end{savequote}

\chapter{\label{ch:1-intro}Introduction} 

\minitoc
\section{Introduction}
Graphs, also called networks, are a natural way of modelling relations in data. Real world examples, such as social, biological and sensor networks exemplify this. However, real-world graph data is often noisy and incomplete -- sensors fail and social network users lie. To handle such issues, the Graph Signal Processing (GSP) community have generalised the tools of sampling and reconstruction from classical signal processing to the irregular graph domain, allowing for a theoretically rigorous foundation; however, the current GSP literature lacks understanding with regards to the impact of varying the amount of information missing, or the impact of the downstream tasks. On the other hand, the Graph Machine Learning (GML) community have approached this issue from a top-down, empirical manner, capturing the full complexity of the real-world tasks under missing data, but without the theoretical rigour.

In this thesis, we focus on two settings. We begin by considering the accuracy of reconstruction of a graph signal with missing node data, focusing on the impact of how many samples we have, which sits well within the GSP literature. We then consider the case where reconstruction of graph signals with missing node data serves some downstream learning task such as classification, and the impact of this task on classification, bridging the GSP and GML literature.

%In this thesis, we focus on the impact of missing data firstly on reconstruction, and then on downstream tasks such as classification.

\section{Motivation}
Incomplete information on graphs occurs in many settings, and in this thesis we focus on cases where the graph is known but node attribute data is missing. We give a non-exhaustive list of motivations for studying how to handle missing noisy node attributes on graphs with practical real-world examples.

\begin{enumerate}
    \item \textbf{Record-keeping}: Data collection at multiple location via sensors (e.g. temperature and wind) can be noisy and incomplete owing to sensor failure, while the geographic location of the sensors (forming the graph) is well known. Reconstruction and denoising allow for filling in of this missing data, and potentially collection of data where failure is likely and sensor maintenance is difficult (e.g. humid or remote locations like rainforests).
    \item \textbf{Cheaper Sensors}: Understanding the trade-offs between sensor failure, noise, and error in the final data collection after denoising and reconstructing missing data can allow for use of cheaper, more numerous sensors and more granular data collection if these downsides of cheaper sensors around noise-sensitivity and failure can be quantified \cite{bush2023impact}.
    \item \textbf{Removal of untrustworthy or private data sources}:  Privacy is important in cases such as healthcare, and untrustworthiness may arise for many reasons (e.g. changes in government have historically lead to lack of trust in particular economic figures in certain countries); removal of data for privacy or trust reasons can leave us with partial data but a complete network. In these cases we might like to replace the removed data with inferred data.
    \item \textbf{Decision making with noisy and incomplete network data}: Decisions are frequently made from noisy and incomplete sensor data -- for example, prediction of solar power input is important in power grid regulation \cite{simeunovic2021spatio}. In this case, we may not care about how accurately we infer the missing data, but rather about minimising the change in our resulting decisions from having incomplete information.

\end{enumerate}
TURN INTO PROSE:
    Equity Trading We give a practical example of these challenges. Prediction of equity prices is lucrative, and stock data is frequently noisy and incomplete (for example, different local markets trade at different hours, causing missing price data - the London Stock Exchange is open well before the New York Stock Exchange). Stocks reflect actual companies with supply chains and established relationships that can be represented as a network. Inference and learning in this setting is thus very valuable, and is naturally suited to an interpretation as a graph learning task with missing feature information.

These listed applications motivate this thesis on incomplete information on graphs. We will advance the understanding of the impact of how much data is missing under noise in different scenarios, and then expand our work to encompass understanding how downstream tasks affect our handling of missing data.

\section{Outline \& Contribution}
The objective of this thesis is to better understand the impact of incomplete data on learning on graphs. We present three technical chapters, based on peer-reviewed or submitted papers, to this end:

\begin{itemize}
    \item \emph{Chapter \ref{ch:proj1_LS}, \nameref{ch:proj1_LS}}:
    This Chapter investigates how sample size affects reconstruction error in the presence of noise via an in-depth theoretical analysis of least-squares reconstruction (LS). Our theorems show that at sufficiently low signal-to-noise ratios (SNRs), under these LS we may simultaneously decrease sample size and decrease average reconstruction error. We further show that at sufficiently low SNRs, for LS reconstruction we have a $\Lambda$-shaped error curve. We also prove that changing sample choice, one of the most common interventions for noise in the literature, cannot mitigate non-monotonicity under LS under mild optimality assumptions. We present thresholds on the SNRs, $\tau$, below which the error is non-monotonic, and illustrate these theoretical results with experiments across multiple random graph models, sampling schemes and SNRs. We prove these results are robust to changes in the signal model, and prove our results for both full-band and bandlimited noise. These results demonstrate that any decision in sample-size choice has to be made in light of the noise levels in the data. 
    \item \emph{Chapter \ref{ch:proj1_GLR}, \nameref{ch:proj1_GLR}}: We continue to investigate how sample size affects reconstruction error in the presence of noise, by extending our in-depth theoretical analysis to cover a regularised reconstruction method, graph Laplacian-regularised reconstruction (GLR). We show that at sufficiently low SNRs, even under a regularised reconstruction method such as GLR we may simultaneously decrease sample size and decrease average reconstruction error. We show that changing sample choice cannot mitigate non-monotonicity without repeated observations of a vertex on a broad class of graphs -- we empirically show that most graphs sampled from Erdos-Renyi, Stochastic Blockmodel or Barabasi-Albert graph models fall into this class.  We further show that at sufficiently low SNRs, for GLR reconstruction, a sample size of $\mathcal{O}(\sqrt{N})$, where $N$ is the total number of vertices, results in lower reconstruction error than near full observation on said class of graphs. We prove these results are robust to changes in the signal model, and prove our results for both full-band and bandlimited noise. We present thresholds on the SNRs, $\tau_{GLR}$, defined in terms of properties of the graph, below which the error is non-monotonic, and illustrate these theoretical results with experiments across multiple random graph models, sampling schemes and SNRs.
    \item \emph{Chapter \ref{ch:proj2}, \nameref{ch:proj2}}: We investigate graph signal reconstruction and sample selection for classification tasks. We present general theoretical characterisations of classification error applicable to multiple commonly used reconstruction methods, and compare that to the classical reconstruction error. We demonstrate the applicability of our results by using them to derive new optimal sampling methods for linearized graph convolutional networks, and show improvement over other graph signal processing based methods on both synthetic and real-world classification tasks.
\end{itemize}

These chapters are based on the following papers:

\begin{itemize}
    \item \emph{Chapter \ref{ch:proj1_LS}}: \textbf{Sripathmanathan, B.,} Dong, X., and  Bronstein, M., ``On the impact of sample size in reconstructing graph signals.'' -- 2023 International Conference on Sampling Theory and Applications (SampTA). IEEE, 2023.
    \item \emph{Chapters \ref{ch:proj1_LS} \& \ref{ch:proj1_GLR}}: \textbf{Sripathmanathan, B.,} Dong, X., and Bronstein, M., ``On the Impact of Sample Size in Reconstructing Noisy Graph Signals: A Theoretical Characterisation.'' -- Under Review at IEEE Transactions on Signal Processing
    \item \emph{Chapter \ref{ch:proj2}}: \textbf{Sripathmanathan, B.,} Dong, X., and Bronstein, M., ``On the Impact of Downstream Tasks on Sampling and Reconstructing Noisy Graph Signals'' -- Under Review at 2025 IEEE International Workshop on Computational Advances in Multi-Sensor Adaptive Processing (CAMSAP).
\end{itemize}

The ideas, figures and content of the technical chapters originate predominantly from me. Furthermore, I contributed to the following paper not presented in the thesis:

\begin{itemize}
    \item Liang, H. *, Borde, HSO. *, \textbf{Sripathmanathan, B. *\footnote{ *Equal Contribution}}, Bronstein, M., Dong, X., ``Towards Quantifying Long-Range Interactions in Graph Machine Learning: a Large Graph Dataset and a Measurement'' -- Under Review at Advances in Neural Information Processing Systems 39 (2025)
\end{itemize}


%%%%%%%%%%%%%%%
%%%% ABSTRACTS:
%%%%%%%%%%%%%%%%
%%%% JOURNAL:

\iffalse
\begin{abstract}
Reconstructing a signal on a graph from noisy observations of a subset of the vertices is a fundamental problem in the field of graph signal processing. This paper investigates how sample size affects reconstruction error in the presence of noise via an in-depth theoretical analysis of the two most common reconstruction methods in the literature, least-squares reconstruction (LS) and graph-Laplacian regularised reconstruction (GLR). Our theorems show that at sufficiently low signal-to-noise ratios (SNRs), under these reconstruction methods we may simultaneously decrease sample size and decrease average reconstruction error. We further show that at sufficiently low SNRs, for LS reconstruction we have a $\Lambda$-shaped error curve and for GLR reconstruction, a sample size of $\mathcal{O}(\sqrt{N})$, where $N$ is the total number of vertices, results in lower reconstruction error than near full observation. We present thresholds on the SNRs, $\tau$ and $\tau_{GLR}$, below which the error is non-monotonic, and illustrate these theoretical results with experiments across multiple random graph models, sampling schemes and SNRs. These results demonstrate that any decision in sample-size choice has to be made in light of the noise levels in the data.
\end{abstract}
\fi

%%%%%% CAMSAP paper:
\iffalse
\begin{abstract}
We investigate graph signal reconstruction and sample selection for classification tasks. We present general theoretical characterisations of classification error applicable to multiple commonly used reconstruction methods, and compare that to the classical reconstruction error. We demonstrate the applicability of our results by using them to derive new optimal sampling methods for linearized graph convolutional networks, and show improvement over other graph signal processing based methods.
\end{abstract}
\fi

%%%%% ATTEMPTS at Motivation
%%% Ch1
\iffalse
In this chapter we investigate the relationship between sample size and reconstruction error, with an emphasis on the phenomenon of reducing the number of observed nodes can reduce reconstruction error and how this phenomenon is mediated by noise. We provide theorems characterising the ways it is possible for this phenomenon to occur for any linear reconstruction method, and then specialise to the case of Least Squares reconstruction to show that the phenomenon does occur. We prove that this phenomenon cannot be mitigated by sample choice if we require a commonly desired form of optimality in choosing said samples, and we demonstrate that this phenomenon continues to happen as the size of the graph gets arbitrarily large. To demonstrate the robustness of our results to changes in the signal and the noise model, we prove a proposition generalising the signal model to any signal model, and we prove equivalent results under bandlimited noise. Finally, we demonstrate this result on synthetic and real-world graph models.
\fi